\documentclass[10pt,a4paper]{amsart}
\usepackage[margin=19mm]{geometry}
\usepackage{amsmath,amssymb,mathtools}
\usepackage[scr=rsfs,cal=euler]{mathalfa}
\usepackage{xfrac}
\usepackage[unicode,hidelinks,bookmarks=false]{hyperref}
\newcommand{\ie}{i.e.\ }
\newcommand{\cf}{cf.\ }
\newcommand{\resp}{respectively\ }
\newcommand{\Subsection}{\subsection}
\providecommand{\nobreakdash}{\nobreak\hbox{-}}
\setlength{\emergencystretch}{2em}
\multlinegap0pt
\usepackage{amssymb}
\usepackage{mathtools}
\usepackage{dsfont}
\usepackage{enumerate}
\usepackage[shortlabels]{enumitem}
\usepackage{float}
\usepackage{tikz}
\newtheorem{proposition}{Proposition}
\newtheorem{remark}{Remark}
\newtheorem{theorem}{Theorem}
\newtheorem{corollary}{Corollary}
  \newcommand{\B}{{\mathcal{B}}}
  \newcommand{\C}{{\mathcal{C}}}
  \newcommand{\G}{{\mathcal{G}}}
\newcommand{\ca}{\mathrm{C}^*}
\newcommand{\foral}{\text{ for all }}
\def\io{\iota}
\def\ka{\kappa}
\newcommand{\wt}{\widetilde}
\title{On hyperrigidity and non-degenerate C*-correspondences}
\author{Joseph A. Dessi, Evgenios T.A. Kakariadis,, Ioannis Apollon Paraskevas}
\date{Source: arXiv:2503.16618v3 (CC BY 4.0)}
\begin{document}
\maketitle

\noindent\emph{Source and licence.} On hyperrigidity and non-degenerate C*-correspondences. Version arXiv:2503.16618v3 (CC BY 4.0), distributed by arXiv under the Creative Commons Attribution 4.0 International licence, http://creativecommons.org/licenses/by/4.0/, as recorded in the arXiv licence metadata for this exact version and shown on the abstract page https://arxiv.org/abs/2503.16618v3. Full licence notice: \texttt{licenses/arxiv-2503.16618-CC-BY-4.0.txt}.\par\smallskip

\noindent\textit{Editorial scope.} Counted unit: the article's theorem T:hypiffuep (source0.tex lines 1646--1659). The article's complete original proof of it is delivered with it (source0.tex lines 1661--1751, one \texttt{proof} environment of 91 lines). No mathematics is written, repaired or re-derived: the statement and the proof are the article's own lines, in the article's own notation, and no retained line is substituted or re-flowed.  Retained uncounted as the source-local context the proof needs: lines 844--847; lines 1034--1037; lines 1190--1198; lines 1259--1261; lines 1326--1335; lines 1346--1349.  Nothing was omitted from any retained span, so the package carries no omission record.  No retained line contains a citation, so the wrapper carries no bibliography and no bibliography entry.  No retained line contains a figure environment, an image-inclusion command or drawing code, so no image is delivered and the package declares no asset dependency.  Editorial formatting only, in addition: an AMS \texttt{amsart} preamble that restates the article's own declarations and macros used by the retained lines, namely the environments \texttt{proposition}, \texttt{remark}, \texttt{theorem}, \texttt{corollary} on separate counters, together with the standard packages those lines need.  The retained environments print in this document as Proposition 1, Remark 1, Theorem 1, Corollary 1 in order of appearance, whereas the article prints the same items with the numbering of its own full document.  The complete native source of the article as distributed in the arXiv e-print for this exact version is bundled unchanged as \texttt{sources/175131/source0.tex}.
\begin{proposition}\label{P:Dil}
Let $\C$ be a C*-algebra and $\phi \colon \C \to \B(H)$ be a completely contractive and completely positive map.
Then there exists a Hilbert space $H' \supseteq H$ and a $*$-homomorphism $\Phi \colon \C \to \B(H')$ such that $\phi = P_H \Phi |_H$.
\end{proposition}

\begin{proposition}\label{P:hypuep}
Let $\G \subseteq \ca(\G)$ be a generating set.
If $\G$ is ucp hyperrigid (resp.\ ccp hyperrigid), then every non-degenerate $*$-representation of $\ca(\G)$ has the ucp unique extension property (resp.\ ccp unique extension property) with respect to $\G$.
\end{proposition}

\begin{proposition}\label{P:uepiffmax}
Let $\G\subseteq \ca(\G)$ be a generating set and $\Phi \colon \ca(\G) \to \B(H)$ be a non-degenerate $*$-representation. 
The following are equivalent:
\begin{enumerate}
\item $\Phi$ has the ccp unique extension property with respect to $\G$.
\item $\Phi$ has the ucp unique extension property with respect to $\G$.
\item $\Phi$ is maximal on $\G$.
\end{enumerate}
\end{proposition}

\begin{remark}\label{R:uepiffmax}
We note for future reference that the proof of the implication [(iii) $\Rightarrow$ (i)] does not require $\Phi$ to be non-degenerate.
\end{remark}

\begin{theorem}\label{T:hypiffmaxun}
Let $\G \subseteq \ca(\G)$ be a generating set such that $\ca(\G)$ is unital.
The following are equivalent:
\begin{enumerate}
\item $\G$ is ccp hyperrigid.
\item $\G$ is ucp hyperrigid.
\item Every non-degenerate $*$-representation $\Phi \colon \ca(\G) \to \B(H)$ is maximal on $\G$.
\end{enumerate}
Moreover we have that $\G$ is ccp/ucp hyperrigid if and only if $\G \cup \{1\}$ is ccp/ucp hyperrigid.
\end{theorem}

\begin{corollary}\label{C:ccp/ucp}
Let $\G\subseteq\ca(\G)$ be a generating set.
Then $\G$ is ccp hyperrigid if and only if $\G$ is ucp hyperrigid.
\end{corollary}

\begin{theorem}\label{T:hypiffuep}
Let $\G \subseteq \ca(\G)$ be a separating generating set.
The following are equivalent:
\begin{enumerate}
\item $\G$ is ccp hyperrigid.
\item $\G$ is ucp hyperrigid.
\item Every non-degenerate $*$-representation $\Phi \colon \ca(\G) \to \B(H)$ has the ccp unique extension property with respect to $\G$.
\item Every non-degenerate $*$-representation $\Phi \colon \ca(\G) \to \B(H)$ has the ucp unique extension property with respect to $\G$.
\item Every non-degenerate $*$-representation $\Phi \colon \ca(\G) \to \B(H)$ is maximal on $\G$.
\item Every $*$-representation $\Phi \colon \ca(\G) \to \B(H)$ is maximal on $\G$.
\item Every $*$-representation $\Phi \colon \ca(\G) \to \B(H)$ has the ccp unique extension property with respect to $\G$.
\item Every $*$-representation $\Phi \colon \ca(\G) \to \B(H)$ has the ucp unique extension property with respect to $\G$.
\end{enumerate}
\end{theorem}

\begin{proof}
The equivalence [(i) $\Leftrightarrow$ (ii)] is Corollary \ref{C:ccp/ucp}, and the implication [(i) $\Rightarrow$ (iii)] is Proposition \ref{P:hypuep}. The equivalence [(iii) $\Leftrightarrow$ (iv) $\Leftrightarrow$ (v)] is Proposition \ref{P:uepiffmax}.
The implication [(vi) $\Rightarrow$ (vii)] is Remark \ref{R:uepiffmax} and the implications [(vii) $\Rightarrow$ (viii) $\Rightarrow$ (iii)] are immediate.
We will complete the proof by showing [(v) $\Rightarrow$ (vi)] and [(vi) $\Rightarrow$ (ii)].

We first prove the implication [(v) $\Rightarrow$ (vi)]. Let $\Phi\colon \ca(\G) \to \B(H)$ and $\Phi'\colon \ca(\G)\to\B(H')$ be $*$-representations such that $H'\supseteq H$ and $\Phi(g)=P_H\Phi'(g)|_H$ for every $g\in\G$.
Moreover, set $L:=[\Phi(\ca(\G))H]$ and $\Phi_0:=\Phi|_L$.
Since $\Phi_0$ is maximal and
\[
P_L \Phi'(g)|_L = P_L \Phi(g)|_L = \Phi_0(g) \foral g \in \G,
\]
we have
\[
\Phi' = \Phi_0 \oplus P_{H' \ominus L} \Phi' |_{H' \ominus L}.
\]
Since
\[
P_{H \ominus L} \Phi'(g) |_{H \ominus L}
=
P_{H \ominus L} \Phi(g) |_{H \ominus L}
=
0
\foral g \in \G,
\]
by the separating property we have $P_{H \ominus L} \Phi' |_{H \ominus L} = 0$.
By positivity of the map $P_{H' \ominus L} \Phi' |_{H' \ominus L}$ and since $\ca(\G)$ is linearly spanned by positive elements we have
\[
P_{H' \ominus L} \Phi' |_{H' \ominus L} = 0_{H \ominus L} \oplus P_{H' \ominus H} \Phi' |_{H' \ominus H}.
\]
Therefore we obtain
\[
\Phi'
= \Phi_0 \oplus P_{H' \ominus L} \Phi' |_{H' \ominus L}
= \Phi_0 \oplus 0_{H \ominus L} \oplus P_{H' \ominus H} \Phi' |_{H' \ominus H}
= \Phi \oplus P_{H' \ominus H} \Phi' |_{H' \ominus H},
\]
i.e., $\Phi$ is a direct summand of $\Phi'$, as required.

We now prove the implication [(vi) $\Rightarrow$ (ii)]. By Theorem \ref{T:hypiffmaxun} it suffices to consider $\ca(\G)$ to be non-unital.
Let $\Phi \colon \ca(\G) \to \B(H)$ be a faithful non-degenerate $*$-representation and let $\phi_n \colon \B(H) \to \B(H)$ be unital completely positive maps such that
\[
\lim_n \|\phi_n(\Phi(g)) - \Phi(g)\| = 0 \foral g \in \G.
\]
Consider the inclusion map
\[
\io \colon \B(H) \to \ell^{\infty}(\B(H)); x \mapsto (x, x, \dots)
\]
and the canonical quotient map
\[
q \colon \ell^{\infty}(\B(H)) \to \ell^\infty(\B(H))/c_0(\B(H)).
\]
Let us also fix a faithful unital $*$-representation
\[
\ka \colon \ell^\infty(\B(H))/c_0(\B(H)) \to \B(K).
\]
Define the faithful $*$-representation
\[
\wt{\Phi} \colon \ca(\G) \to \B(K); \wt{\Phi}(c) := \ka \circ q \circ \io \circ \Phi(c),
\]
and the completely contractive completely positive map
\[
\wt{\Psi} \colon \ca(\G) \to \B(K); \wt{\Psi}(c) = \ka \circ q ( \phi_1(\Phi(c)), \phi_2(\Phi(c)), \dots ).
\]
Proposition \ref{P:Dil} implies that there exists a $*$-representation $\wt{\Phi}'\colon \ca(\G)\to \B(K')$ such that $K'\supseteq K$ and
\[
\wt{\Psi}(c)=P_K \wt{\Phi}'(c)|_K \foral c\in \ca(\G).
\]
For every $g \in \G$ we have
\[
P_K \wt{\Phi}'(g)|_K
=
\wt{\Psi}(g)
=
\ka \circ q ( \phi_1(\Phi(g)), \phi_2(\Phi(g)), \dots )
=
\ka \circ q \circ \io \circ \Phi(g)
=
\wt{\Phi}(g).
\]
By our assumption $\wt{\Phi}$ is maximal on $\G$, and therefore we obtain that $K$ is reducing for $\wt{\Phi}'$.
In particular we have $\wt{\Psi}(c)=\wt{\Phi}(c)$ for every $c\in\ca(\G)$.
Equivalently, we have
\[
\ka \circ q(\phi_1(\Phi(c)), \phi_2(\Phi(c)), \dots)
=
\ka \circ q(\Phi(c), \Phi(c), \dots)
\foral
c \in \ca(\G).
\]
Since $\ka$ is faithful we conclude that the equality $\lim_n \phi_n(\Phi(c)) = \Phi(c)$ holds for all $c \in \ca(\G)$, as required.
\end{proof}

\end{document}
